O gerenciamento lógico da infraestrutura de rede compreende o planejamento, implementação e manutenção da configuração da camada 2. Esta atividade é estratégica para garantir isolamento adequado entre serviços e desempenho da rede. As atividades constantes relacionadas a esta responsabilidade são: 

\begin{itemize} 
    \item Planejamento e criação de VLANs para segmentação lógica da rede; 
    \item Configurações de swicthes;
    \item Apoio na resposta a incidentes de segurança relacionados à rede; 
\end{itemize}


    \item Implementação de NAT quando necessário; 
    \item Monitoramento de tráfego para identificação de gargalos ou comportamentos anômalos; 
    \item Revisão periódica das regras de firewall e políticas de acesso; 
    \item Definição de políticas de comunicação entre VLANs (no HP são as Zonas); 
    \item Criação e manutenção de regras de firewall para controle de acesso entre redes internas e externas; 
    \item Planejamento de endereçamento IP e organização de sub-redes; 